\documentclass[a4paper,fleqn]{cas-sc}
%% Swap to cas-dc for the double-column layout: \documentclass[a4paper,fleqn]{cas-dc}
\usepackage[numbers]{natbib}

\usepackage{amssymb}
\usepackage{amsmath}
\usepackage{amsthm}
\usepackage{orcidlink}

\newtheorem{theorem}{Theorem}[section]
\newtheorem{lemma}[theorem]{Lemma}
\newtheorem{proposition}[theorem]{Proposition}
\newtheorem{corollary}[theorem]{Corollary}
\theoremstyle{definition}
\newtheorem{definition}[theorem]{Definition}
\newtheorem{example}[theorem]{Example}
\theoremstyle{remark}
\newtheorem{remark}[theorem]{Remark}

\newcommand{\Lam}{\mathrm{Lam}}
\newcommand{\LamDB}{\mathrm{Lam}_{\mathrm{db}}}
\newcommand{\alphaeq}{=_{\alpha}}


\begin{document}

\shorttitle{A Formalized Foundation for Research in Untyped Lambda Calculus}
\shortauthors{Chris Anto Fröschl et al.}

\title[mode=title]{A Formalized Foundation for Research in Untyped Lambda Calculus}

\author[1]{Chris Anto Fröschl}[orcid=0009-0006-3562-5395]
\cormark[1]
\ead{chris.froeschl@hm.edu}

\author[1]{Prof. Matthias Güdemann}[orcid=0000-0002-1002-6023]
\cormark[1]
\ead{matthias.guedemann@hm.edu}

\affiliation[1]{
    organization={Munich University of Applied Sciences, Department of Computer Science and Mathematics},
    addressline={Lothstraße 34},
    city={Munich},
    postcode={D-80335},
    state={Bavaria},
    country={Germany}
}

\cortext[1]{Corresponding author}

\begin{abstract}
TODO Intro CSLib, Crole 5 equivs, dB ~= QAEQUIV and what it builds on
Derived results from it, "Formalized Barendregt"
Perhaps even benchmark like lambda n ways?

Lorem ipsum dolor sit amet, consectetur adipiscing elit, sed do eiusmod tempor incididunt ut labore et dolore magna aliqua. Ut enim ad minim veniam, quis nostrud exercitation ullamco laboris nisi ut aliquip ex ea commodo consequat. Duis aute irure dolor in reprehenderit in voluptate velit esse cillum dolore eu fugiat nulla pariatur. Excepteur sint occaecat cupidatat non proident, sunt in culpa qui officia deserunt mollit anim id est laborum.
Lorem ipsum dolor sit amet, consectetur adipiscing elit, sed do eiusmod tempor incididunt ut labore et dolore magna aliqua. Ut enim ad minim veniam, quis nostrud exercitation ullamco laboris nisi ut aliquip ex ea commodo consequat. Duis aute irure dolor in reprehenderit in voluptate velit esse cillum dolore eu fugiat nulla pariatur. Excepteur sint occaecat cupidatat non proident, sunt in culpa qui officia deserunt mollit anim id est laborum.
\end{abstract}

\begin{highlights}
\item First machine-checked formalization of Crole's five equivalent
  definitions of alpha equivalence.
\item An isomorphism, formalized in Lean~4, between a de Bruijn
  representation and the quotient of named terms by alpha
  equivalence, following Norrish's proof pearl.
\item A roadmap toward a formalized foundation for untyped lambda
  calculus research, including further binding representations and
  a "formalized Barendregt".
\end{highlights}

\begin{keywords}
$\alpha$-equivalence \sep
Atom \sep
$\lambda$-expression \sep
Permutation action \sep
Renaming \sep
Variable binding \sep
de Bruijn indices \sep
formal verification \sep
proof assistant \sep
Lean \sep
\MSC[2020]{03B70 \sep 68V15 \sep 68V20}
\end{keywords}

\maketitle

\section{Introduction}
\label{sec:introduction}

Variable binding is the recurring technical obstacle in any
formalization touching the lambda calculus: essentially every
representation choice (named terms modulo alpha equivalence,
de Bruijn indices, locally nameless terms, higher-order abstract
syntax, nominal sets, \ldots) trades one set of proof obligations for
another, and results proved for one representation rarely transfer
automatically to the others. This has led to a large but fragmented
body of formalization work, often redone from scratch for each new
proof assistant or each new project, with little of it packaged as a
reusable library.

CSLib is an effort to build such a reusable, foundational library for
computer science results in the Lean~4 theorem prover, in the spirit
of Mathlib for mathematics~\cite{demoura2026a}. This paper reports on
two contributions to CSLib's treatment of the untyped lambda calculus:

\begin{enumerate}
  \item a machine-checked formalization of Crole's result that five
    common literature definitions of alpha equivalence on named lambda terms
    are equivalent~\cite{crole2012a}, which to our knowledge
    had not previously been formalized in a proof assistant; and
  \item an isomorphism, following Norrish's proof
    pearl~\cite{norrish2007a}, between the quotient of named
    terms by a chosen alpha equivalence and a de Bruijn
    representation, connecting the named and nameless ends of
    CSLib's binding infrastructure.
\end{enumerate}

We present both results as part of a longer-term plan: to grow CSLib
into an idiomatic, well-reviewed, and maintained foundation that
other researchers can build on directly, rather than reformalizing
lambda calculus basics for each new project. Section~\ref{sec:future-work}
sketches this plan in more detail, including further equivalences in
the style of Crole's for SK combinator calculus and other binding
representations, and a longer-term goal of formalizing a substantial
part of Barendregt's reference treatment of the untyped lambda
calculus~\cite{barendregt1985a}.

\paragraph{Outline.}
Section~\ref{sec:cslib} introduces CSLib and its current named
representation of the untyped lambda calculus.
Section~\ref{sec:alpha-equivalences} presents the formalization of
Crole's alpha-equivalence equalities.
Section~\ref{sec:de-bruijn} presents the isomorphism with the de
Bruijn representation. Section~\ref{sec:related-work} discusses
related formalization efforts, Section~\ref{sec:future-work} lays
out planned extensions, and Section~\ref{sec:conclusion} concludes.

\section{CSLib and its named representation of the lambda calculus}
\label{sec:cslib}

\subsection{CSLib}
\label{subsec:cslib-overview}

CSLib is a community-maintained library of foundational computer
science formalizations for the Lean~4 theorem prover, built on top
of Mathlib. \cite{demoura2026a} 
Its goal is to provide idiomatic, well-reviewed, and actively
maintained formalizations of standard results, so that later
research projects can build directly on a shared, trusted base
rather than re-deriving basic definitions and lemmas from scratch.
The work reported in this paper concerns CSLib's treatment of the
untyped lambda calculus, and in particular its representation(s) of
variable binding.

\subsection{Named terms}
\label{subsec:named-terms}

CSLib currently represents untyped lambda terms in a standard named
style: variables are drawn from a countably infinite set of atoms,
and terms are built from variables, application, and abstraction. TODO dummy math

\begin{definition}[Named lambda terms]
\label{def:named-terms}
Let $\mathrm{Var}$ be a countably infinite set of variables. The set
$\Lam$ of named lambda terms is generated by
\[
  t \;::=\; x \mid t\, t \mid \lambda x.\, t
  \qquad (x \in \mathrm{Var}).
\]
\end{definition}

On top of this definition, CSLib provides the usual auxiliary
notions needed to state and prove properties of the calculus: free
variables, capture-avoiding substitution, and a baseline notion of
alpha equivalence used as the reference point against which
Crole's alternative definitions are shown equivalent in
Section~\ref{sec:alpha-equivalences}.


\section{Formalizing Crole's alpha-equivalence equalities}
\label{sec:alpha-equivalences}

Crole~\cite{crole2012a} gives five definitions of alpha
equivalence on named lambda terms, arising from different intuitions
about what it means to rename bound variables (via explicit
substitution, via swapping, via simultaneous renaming of bound
variable occurrences, \ldots), and proves that all five coincide.
The proof is unremarkable to state on paper, but formalizing it
exposes a number of routine-but-fiddly inductive arguments about
substitution and renaming under binders that are easy to state
informally and easy to get subtly wrong when made fully precise ---
which is precisely the kind of result CSLib aims to package once and
for all.

\subsection{The five definitions}
\label{subsec:five-definitions}

\subsection{Equivalence of the five definitions}
\label{subsec:equivalence-theorem}

TODO theorems + proof links + sketch

\begin{theorem}[Crole~\cite{crole2012a}, formalized]
\label{thm:crole-equivalence}
The five relations of Section~\ref{subsec:five-definitions} coincide
as relations on $\Lam$.
\end{theorem}

\subsection{Formalization notes}
\label{subsec:formalization-notes}

TODO extra work due to the paper handing to the literature


\section{An isomorphism with de Bruijn terms}
\label{sec:de-bruijn}

TODO: fill in the actual de Bruijn development being reused
(source/authorship, what was already in CSLib/Mathlib vs. what was added),
the precise choice of which of the five equivalent alpha
equivalences is quotiented by, and the isomorphism statement and
proof sketch, following Norrish's proof pearl.

De Bruijn indices remove the need for alpha equivalence altogether by
replacing bound variable names with numerals counting binder depth.
Norrish's proof pearl~\cite{norrish2007a} demonstrates that,
despite a reputation for being awkward to reason about, a de Bruijn
representation supports clean, structurally recursive proofs of the
standard metatheory of the lambda calculus once the representation is
set up carefully. We reuse an existing de Bruijn development
and show that it is isomorphic to the quotient of named terms
$\Lam$ by (any one of) the alpha equivalences of
Section~\ref{sec:alpha-equivalences}.

\subsection{De Bruijn terms}
\label{subsec:db-terms}

\begin{definition}[De Bruijn terms]
\label{def:db-terms}
The set $\LamDB$ of de Bruijn lambda terms is generated by
\[
  t \;::=\; n \mid t\, t \mid \lambda.\, t
  \qquad (n \in \mathbb{N}).
\]
\end{definition}

\subsection{The isomorphism}
\label{subsec:isomorphism}

\begin{theorem}
\label{thm:db-isomorphism}
$\Lam / \alphaeq \;\cong\; \LamDB$.
\end{theorem}


\section{Related work}
\label{sec:related-work}

Formalizing variable binding has a long history across proof
assistants, and several general-purpose approaches exist alongside
the direct, elementary style used in CSLib. Nominal
techniques~\cite{pitts2013a}, as implemented in Nominal
Isabelle~\cite{urban2005a}, provide binder-generic reasoning
principles via permutations and freshness, at the cost of adopting a
dedicated meta-theory. Locally nameless
representations~\cite{aydemir2008a} aim to make binder-heavy formalizations more
tractable or more automatically generated from a specification,
respectively. Higher-order abstract syntax sidesteps binder
representation issues by delegating them to the proof assistant's
own binding mechanism, again at the cost of restricting which proof
techniques remain available (e.g.\ induction principles).

CSLib's current approach is deliberately more elementary: named terms
and de Bruijn terms are both represented directly as first-order
inductive types, with the relationships between representations
(such as the isomorphism of Section~\ref{sec:de-bruijn}) proved
explicitly rather than assumed via a shared meta-level binding
mechanism. This trades some proof automation for representations
that are simple to state, and results that transfer cleanly between
representations by construction, which fits CSLib's goal of a
reusable foundation rather than a self-contained framework.


\section{Future work}
\label{sec:future-work}

TODO sketch general net of equivalences in Proof pearls aswell (image)

The two results of this paper are intended as the first pieces of a
larger, ongoing effort within CSLib to build a formalized foundation
for research on untyped lambda calculus. We outline the planned
directions below.

\subsection{Further equivalence results}
\label{subsec:further-equivalences}

Crole-style results, showing several natural but superficially
different definitions of the same notion to coincide, are not unique
to alpha equivalence. We plan analogous formalizations for SK
combinator calculus and other binding representations (locally
nameless terms, nominal terms), each intended to slot into CSLib
alongside the results of this paper and to be connected to them by
explicit isomorphisms in the style of Section~\ref{sec:de-bruijn}.

\subsection{A formalized Barendregt}
\label{subsec:formalized-barendregt}

Barendregt's \emph{The Lambda Calculus: Its Syntax and
Semantics}~\cite{barendregt1985a} remains the standard reference
for the untyped lambda calculus, and a substantial fraction of its
content --- confluence, standardization, the theory of reduction,
combinatory completeness, and the associated notions of models ---
has never been formalized in a single, coherent development. We see
CSLib's binding infrastructure as a natural basis for such an effort,
which we tentatively refer to as a "formalized Barendregt". Beyond
its intrinsic value as a reusable library, a sufficiently complete
formalized Barendregt would put open or only recently resolved
conjectures discussed in Barendregt's later survey work within reach
of machine-checked treatment.

\subsection{Maintenance and reuse}
\label{subsec:maintenance}

Throughout, the intent is for this work to be maintained as part of
CSLib rather than as a standalone, one-off development: subject to
code review, kept up to date with CSLib and Mathlib, and structured
so that later projects can depend on it directly.

\section{Conclusion}
\label{sec:conclusion}

We have presented CSLib's current named representation of the
untyped lambda calculus, a formalization of Crole's result that five
notions of alpha equivalence coincide, and an isomorphism, following
Norrish, between the quotient of named terms under alpha equivalence
and a de Bruijn representation. Together these results connect two
of the standard approaches to variable binding within a single,
reusable Lean~4 library, and form the starting point for a longer-term
effort to build a maintained, idiomatic foundation for research on
untyped lambda calculus, including further binding representations
and a formalized treatment of Barendregt's reference text.



\section*{Data availability}
The Lean~4 formalization discussed in this paper is developed as
part of CSLib and is publicly available at
\url{https://github.com/leanprover/cslib}.

TODO consider with zenodo

\bibliographystyle{cas-model2-names}
\bibliography{refs}

\end{document}
